\subsection{Inflation and a restricted rank-eight endpoint reduction}
\label{subsec:k12-coloop-rank8}

The orthogonal-sum operation used for the explicit coface is the
inflation construction of \cite[\S6.1]{ElbazGanglSoule2013}. We record
its chain-level consequence and its limits before giving a further
deduction for the rank-eight endpoint.

\begin{lemma}[Inflation, with the incidence convention fixed]
\label{lem:k12-inflation}
Let $C_d(N;R)$ be the oriented cellular coinvariants of the relative
projective Voronoi complex. Inflation defines a degree-one map
\[
 J_N:C_d(N;R)\longrightarrow C_{d+1}(N+1;R),\qquad
 J_N(\tau)=\tau\perp\mathbb Z.
\]
With the cone-span orientation convention of
\cite[\S1.2.1, equations (1)--(2)]{DEKM2026}, and the new quadratic
ray appended to an oriented basis, one has
$\partial J_N=-J_N\partial$. Thus an explicit equation
$\partial b=\tau$ gives $\partial(-J_Nb)=J_N\tau$, with coefficient
one. If two is invertible in $R$, then $J_{N+1}J_N=0$.
\end{lemma}

\begin{proof}
Normalize a form witnessing the minimum set of $\tau$ to minimum one.
Its orthogonal sum with $[1]$ has exactly the old minimum pairs and the
new coordinate pair: a vector with both components nonzero has norm
at least two. The resulting projective cell is the cone on $\tau$.
For a facet $\rho$ of $\tau$, the basis obtained by appending an old
missing ray to the oriented basis of $J_N\rho$ differs by one
transposition from the corresponding basis for $J_N\tau$.
Its incidence is therefore the negative of the old incidence.
The additional base facet has vector rank $N$ in the rank-$(N+1)$
ambient lattice and is zero in the relative complex. Old relative
boundary faces remain relative boundary faces after inflation.
This proves the differential identity, also for nonsimplicial cells.

The construction precedes passage to coinvariants. An old integral
equivalence $g$ extends to $g\perp1$, and an old orientation reverser
extends to one for the inflated cell. Thus the map descends, with no
stabilizer-index multiplier; actual equivalent facets are summed with
their ordinary incidence signs on both sides. Finally the integral
permutation of the two added coordinate axes fixes the old quadratic
span and transposes two independent new rays. It reverses orientation,
proving the double-inflation cancellation after inverting two.
\end{proof}

A rational coloop is a minimum ray whose removal lowers the vector
rank. It is \emph{integrally split} if, for its primitive representative
$v$ and the span $W$ of the other minima,
$L=(L\cap W)\oplus\mathbb Zv$. Such a cell is an inflation, even when
the original witness form has nonzero cross terms. If the minima
generate $L$, every rational coloop splits integrally. In general the
glue in an overlattice can prevent this.

Lemma~\ref{lem:k12-inflation} is not a contracting homotopy for the
subcomplex spanned by inflated cells. A cycle in the image of $J_N$
need not lift to a cycle, since its lifted boundary can lie in
$\ker J_N$. Consequently lower-rank homology vanishing alone does not
eliminate all cycles supported on such cells. The known small-rank
splitting in \cite[Theorem~6.1]{ElbazGanglSoule2013} is a separate
statement; \cite[\S6.3]{ElbazGanglSoule2013} also describes the failure
of a splitting at the next rank.

\begin{proposition}[K17-K12-EXCESS5]
\label{prop:k12-excess-five}
Let $L$ be a well-rounded lattice of rank $N\geq6$ whose perfect cone
$\sigma$ is nonsimplicial and has perfection rank $N+5$.
There is a rational subspace $W$ of dimension six such that the
minimum rays in $W$ span a nonsimplicial cone of dimension eleven.
Outside $W$ there are exactly $N-6$ minimum pairs, whose images form
a basis modulo $W$; each is a rational coloop.

If $N\geq8$ and the minima generate $L$, the cone has an integral
orientation-reversing stabilizer. In particular, a nonsimplicial
rank-eight cell of projective degree twelve with minimum-generated
ambient lattice is zero in oriented coinvariants over
$R=\mathbb Z[1/13!]$.
\end{proposition}

\begin{proof}
Choose a circuit of $k$ minimum quadratic rays and let their vector
span be $W$, of dimension $m$. Its quadratic rank is $k-1$.
The positive and negative parts of a circuit relation are equal
positive semidefinite sums. Comparing kernels shows that both
supports span $W$, so $k\geq2m$.

The cone $\sigma_W$ of all minima in $W$ is a perfect cone for
$L\cap W$ and is a face of $\sigma$: a positive semidefinite
functional with kernel $W$ exposes it. It is nonsimplicial, since
distinct rank-one positive semidefinite generators are extremal and
the selected circuit lies among them. Write $r_W=\dim\sigma_W$.
Choosing $N-m$ minima independent modulo $W$ and projecting their
quadratic rays gives
\[
 N+5\geq r_W+N-m,
 \qquad r_W\geq k-1\geq2m-1.
\]
Thus $m\leq6$ and $r_W\leq m+5\leq11$.
By \cite[Theorem~1.1(ii)]{DutourHulekSchurmann2015}, the only
nonsimplicial perfect cone of dimension at most ten is the $D_4$
cone. That exception would have $m=4$ and $r_W=10$, contradicting
$r_W\leq m+5$. Hence $r_W=11$ and $m=6$. The chosen circuit has
exactly twelve rays.

Choose minima $v_1,\ldots,v_t$ independent modulo $W$, where
$t=N-6$. Since $\dim\sigma=11+t$, their quadratic rays form a
basis modulo $\operatorname{span}(\sigma_W)$. Choose a vector-space
complement $U$ and use their images as its basis. For another
minimum vector $v\notin W$, express $vv^t$ in this quadratic basis,
modulo a form supported on $W$. If the $U$-coordinates of $v$ are
$a_1,\ldots,a_t$, the $U\times U$ block is diagonal on the right.
Therefore $a_i a_j=0$ for $i\ne j$, so exactly one coordinate,
say $a_j=a$, is nonzero. The coefficient of $v_jv_j^t$ is $a^2$.
The $W\times U$ block then shows that the $W$-component of $v$
is $a$ times that of $v_j$. Thus $v=av_j$, and equal minimum norms
give $a=\pm1$. There are precisely the stated $t$ outside pairs.

If the minima generate $L$, let $M$ be generated by those in $W$.
The complete minimum list gives
\[
 L=M\oplus\bigoplus_{j=1}^t\mathbb Zv_j,
 \qquad M=L\cap W.
\]
For $N\geq8$, interchange $v_1,v_2$ and fix $M$ and all remaining
$v_j$. This is an integral lattice automorphism. It fixes the
eleven-dimensional quadratic span of $\sigma_W$ and transposes
two independent outside rays, hence reverses the cell orientation.
The coinvariant therefore vanishes after inverting two.
\end{proof}

This proposition alone leaves simplicial rank-eight cells and nonsimplicial
cells with nontrivial ambient glue. The following theorem removes the
latter as well. Neither result asserts the rank-eight
$H_5$ vanishing, enumerate the complete degree-twelve chain group,
or compute its incoming boundary matrix. The source's low-dimensional
perfect-cone classification is invoked as cited background; its
enumeration was not rerun in this deduction.
